\documentclass{situationist-preprint}

\title{The Existence of Qualia}
\msid{MND-10-a}
\msfamily{MND-10}
\msdiscipline{Philosophy of Mind}
\mstype{Principal manuscript}
\mskeywords{qualia; phenomenal character; illusionism; representationalism; intrinsic properties; quality spaces; structuralism; Newman's objection}

\begin{document}
\maketitle

\begin{abstract}
Qualia are the properties in virtue of which there is something it is like to undergo an experience. Illusionists deny that there are any, holding that experiences only seem to have them; representationalists deny that they are anything over and above what experiences represent; and structuralists hold that a quality is exhausted by its place in a space of similarities. We argue that qualia exist, that they are intrinsic to the experiences that have them, and that their structure can nonetheless be described in purely relational terms. The existence claim rests on the reliability of introspection about current phenomenal states and on a dilemma for illusionism: the seeming that it posits is either itself phenomenal, in which case qualia exist, or purely cognitive, in which case it does not explain the datum. The intrinsicness claim rests on an application of Newman's objection to structuralism: a relational description of a quality space is satisfied by any domain with the right cardinality and structure, so it cannot say which qualities occupy the space. We show that a relational description is complete up to structure-preserving permutation and leaves out precisely the intrinsic character of each quality, which explains why qualia seem ineffable.
\end{abstract}

\section{The problem}

When one looks at a ripe tomato, tastes coffee or stubs a toe, there is something it is like to have the experience, and what it is like differs from case to case \citep{Nagel1974}. \emph{Qualia} is the name philosophers have given to the properties of experiences in virtue of which this is so. The question whether there are qualia is not the question whether there are experiences, which few deny, but whether experiences have properties of this kind, and if so whether these properties can be identified with properties of some other, less puzzling kind.

Three positions dominate the debate. \emph{Realists} hold that qualia exist and are not exhausted by the functional roles or representational contents of the states that have them \citep{Block1990,Block1996,Chalmers1996}. \emph{Eliminativists} and \emph{illusionists} hold that there are no qualia: experiences do not have the properties we think they have, and the task of a theory of consciousness is to explain why we think so \citep{Dennett1988,Frankish2016}. \emph{Representationalists} hold that the phenomenal character of an experience is identical to, or supervenes on, its representational content, so that what it is like to see red is a matter of red's being represented in a certain way \citep{Harman1990,Tye1995}; on this view there are phenomenal properties, but they are not qualia in the realist's sense. A fourth position, which we shall call \emph{structuralism}, holds that each sensory quality is wholly constituted by its relations of similarity and difference to the others, so that a complete description of the quality space is a complete description of the qualities \citep{Clark1993,Lee2021}.

We argue for a view that combines the realist's thesis with the structuralist's insight. Qualia exist and are intrinsic properties of experiences, and their structure is fully describable in relational terms. The two claims are compatible because a relational description, however complete, fixes its domain only up to isomorphism.

\section{A short history}

The term was introduced in its modern sense by \citet{Lewis1929}, for whom qualia were recognizable qualitative characters of the given, repeatable from one experience to another and in themselves ineffable. The verificationist and behaviorist climate of the following decades left little room for them, and when they returned it was in arguments against functionalism. \citet{Shoemaker1982} asked whether two subjects could be functionally alike while their color experiences were inverted, and argued that, although qualia are functionally definable as a class, inversion within the class cannot be ruled out. \citet{Nagel1974} argued that the subjective character of experience is inaccessible to any objective description, and \citet{Jackson1982} that a scientist who knows all the physical facts about color vision but has never seen color learns something on first seeing it. The knowledge argument is answered in MND-04, which holds that Mary learns no new fact but acquires new concepts of facts she already knew; the existence of qualia is not at issue there, since the concepts she acquires are concepts of the qualities she now experiences.

The modern case against qualia was made by \citet{Dennett1988}. Qualia, he argued, are supposed to be ineffable, intrinsic, private and directly apprehensible, and no property has all four features: the intuition pumps that seem to show them, such as the coffee tasters who disagree over whether their experience or their preference has changed, show instead that there is no fact of the matter about which features an experience has. \citet{Dennett1991} extended the argument into a general theory of consciousness. \citet{Harman1990} added the argument from transparency: when one tries to attend to the qualities of one's experience, one finds only the qualities of the objects experienced, so that the alleged intrinsic qualities of experience are not introspectively available. \citet{Tye1995} developed this into a theory on which phenomenal character is a species of representational content. \citet{Block1990,Block1996} replied with cases, such as Inverted Earth, in which representational content and phenomenal character seem to come apart, and with the distinction between the properties an experience represents and the \enquote{mental paint} by which it represents them. \citet{Kind2003} argued that the transparency of experience is at most partial.

The most recent stage has been the development of illusionism into an explicit research programme. \citet{Frankish2016} argues that the realist's notion of phenomenal consciousness is a theoretical posit that should be abandoned, and that what requires explanation is the illusion of phenomenality: the fact that introspection represents experiences as having qualitative properties they lack. \citet{Chalmers2018} formulates the \enquote{meta-problem} of explaining why we judge that there is a hard problem, and argues that illusionism is the view on which the meta-problem is the only problem. Critics have argued that the illusion of phenomenality would itself be phenomenal \citep{Balog2016,Schwitzgebel2016}. We develop that response into a dilemma below.

Meanwhile, work on quality spaces, from \citet{Clark1993} through \citet{Palmer1999} to \citet{Lee2021}, has shown how much of the character of experience can be captured by relational structure: color space, with its hue circle, its unique hues and its asymmetries \citep{Hardin1988}; pitch space, with its helical structure of height and chroma; and the dimensional structure of taste and smell. Structuralists have concluded that nothing else is needed. We accept the descriptive achievement and resist the conclusion.

\section{Framework}

\begin{definition}[Quale]\label{d:quale}
A \emph{quale} is a property of an experience in virtue of which there is something it is like for its subject to undergo it, such that experiences that differ in their qualia differ in what it is like to undergo them.
\end{definition}

Definition~\ref{d:quale} does not build in ineffability, privacy or infallible apprehension. These were part of Lewis's notion and of the notion Dennett attacked, but none of them is needed to formulate the question, and the realist does better to argue for whichever of them is true than to stipulate them.

\begin{definition}[Intrinsic]\label{d:intrinsic}
A property $F$ of an experience $e$ is \emph{intrinsic} to $e$ if any duplicate of $e$, that is, any experience with the same constituents standing in the same relations among themselves, has $F$. A property is \emph{relationally constituted} if whether a thing has it depends on that thing's relations to things that are not its constituents.
\end{definition}

The definition applies Lewis's test of intrinsicness to experiences. It is neutral on what an experience's constituents are. On views of perception on which a veridical perceptual experience has the perceived object as a constituent, as on the naïve realism defended in MND-16, the qualia of that experience can be intrinsic to it even though the experience is a relation between subject and object. Qualia of bodily sensation, imagery, dreams and hallucination are intrinsic to experiences that have no external constituents.

\begin{definition}[Quality space; relational description]\label{d:space}
A \emph{quality space} for a sensory modality is a structure $\langle Q, R_1, \ldots, R_n\rangle$ where $Q$ is the set of qualia of that modality and the $R_i$ are relations among them, such as exact similarity, relative similarity (\enquote{$a$ is more similar to $b$ than to $c$}), betweenness and opponency. A \emph{relational description} of the space is a description that specifies which tuples of members of $Q$ stand in which relations, without naming any member of $Q$ except as the occupant of a place in this structure.
\end{definition}

A relational description is, in effect, the Ramsey sentence of a theory of the modality's qualities \citep{Lewis1970}: it says that there are some qualities standing in such-and-such relations, without saying which they are.

\begin{principle}[Introspective reliability]\label{p:intro}
Introspective judgments about the current phenomenal character of attended experiences, under normal conditions, are reliable.
\end{principle}

Principle~\ref{p:intro} is the conclusion of MND-20, which argues that introspection is reliable about current phenomenal states and unreliable about their causes. We use only the first half.

\begin{principle}[Distinctness of seeming]\label{p:seem}
A state in which it seems to a subject that her experience has a certain character is either itself an experience with phenomenal character, or a cognitive state, such as a belief or an introspective representation, that has no phenomenal character of its own.
\end{principle}

Principle~\ref{p:seem} is a dichotomy, not a substantive thesis: a state either has phenomenal character or it does not.

\section{The existence of qualia}

\begin{lemma}\label{l:datum}
There are experiences that, by reliable introspection, have properties in virtue of which there is something it is like to undergo them, and that differ in what it is like to undergo them.
\end{lemma}

\begin{proof}
Subjects judge, by introspection of current attended experience, that seeing red differs from seeing green, and that pain differs from pressure. These are judgments about current phenomenal character at a grain that subjects are trained to discriminate. By Principle~\ref{p:intro}, most such judgments are true. They are judgments that the experiences have properties in virtue of which there is something it is like to undergo them, and in respect of which they differ.
\end{proof}

\begin{lemma}\label{l:illusion}
Illusionism, the thesis that experiences merely seem to have qualia while having none, is false.
\end{lemma}

\begin{proof}
Illusionism holds that there is a state in which it seems to the subject that her experience has phenomenal character, and that the seeming is non-veridical. By Principle~\ref{p:seem}, this seeming state either has phenomenal character itself or does not. If it does, then there is at least one state with phenomenal character, and that state has a property in virtue of which there is something it is like to be in it; by Definition~\ref{d:quale}, it has qualia, and illusionism is false. If it does not, the seeming is a cognitive representation with no character of its own. Then what the illusionist explains is that subjects represent their experiences as having qualitative character. But the datum of Lemma~\ref{l:datum} is not that subjects so represent them; it is that the experiences, by reliable introspection, have the character. To treat the datum as a mere representation is to deny Principle~\ref{p:intro} for exactly the class of judgments to which it applies, and the illusionist gives no independent reason to do so. Either way, illusionism fails.
\end{proof}

\begin{proposition}\label{p:exist}
Qualia exist.
\end{proposition}

\begin{proof}
By Lemma~\ref{l:datum} there are experiences with properties satisfying Definition~\ref{d:quale}, and by Lemma~\ref{l:illusion} the illusionist's alternative interpretation of the datum fails.
\end{proof}

The second horn of the dilemma is where illusionists will press, and we return to it in Objection~\ref{o:frankish}. Here we note only that the strength of the argument lies in its asymmetry. The illusionist needs introspection to be systematically wrong about the one thing it is best placed to get right, and the evidence of introspective error, reviewed in MND-20, concerns causes, memory and theory-laden classification, not the presence of current phenomenal character.

\section{Intrinsicness and relational structure}

We now argue that qualia are not relationally constituted, either by their place in a quality space or by their relations to what experiences represent.

\begin{lemma}[Newman's point]\label{l:newman}
A relational description of a quality space is satisfied by any domain of the right cardinality on which relations with the specified extensions can be defined. In particular, if $\pi$ is a permutation of $Q$ that preserves every $R_i$, then the relational description is satisfied equally when each quality $q$ occupies the place it occupies and when each $q$ occupies the place of $\pi(q)$.
\end{lemma}

\begin{proof}
The first claim is the point made by \citet{Newman1928} against Russell's structural realism: a statement that there exist entities with a given structure, together with a specification of their number, is true of any collection of that number, since relations with any desired extension can be defined on it. The second claim follows from the definition of a relational description (Definition~\ref{d:space}): such a description names no quality except as the occupant of a place, and a structure-preserving permutation maps the structure onto itself, so every sentence of the description has the same truth value under the permuted assignment.
\end{proof}

\begin{lemma}\label{l:notrel}
Qualia are not constituted by their places in a quality space.
\end{lemma}

\begin{proof}
Suppose that each quale were constituted by its place in the quality space. Then, by Lemma~\ref{l:newman}, which quale occupies which place would not be a further fact, and the experience of any domain of the right structure would be the experience of the same qualities. But the domain of auditory pitches and a suitably restricted domain of color hues can be given isomorphic structures for some relations; more simply, a domain of numbers can be given any structure we like, and it does not thereby become a domain of colors. What makes a domain a domain of colors, rather than pitches or numbers, is what its members are like, and this is not fixed by the relations among them. So the qualia are not constituted by their places: they occupy them in virtue of what they are.
\end{proof}

\begin{lemma}\label{l:notrep}
Qualia are not constituted by the relations of experiences to the properties they represent.
\end{lemma}

\begin{proof}
If the qualia of an experience were constituted by what it represents, then any two experiences alike in representational content would be alike in qualia. In cases of the kind described by \citet{Block1990}, a subject who moves to Inverted Earth, where the colors of things are inverted and her lenses invert them back, comes after long immersion to represent the sky as blue, as the natives do, while what it is like for her to see the sky remains what it was. More simply, an experience with no external objects, such as an afterimage or a pain, has qualia whose character does not vary with what it represents, since its representational relations to the environment are minimal or absent. Hence the relation to represented properties does not constitute qualia.
\end{proof}

\begin{theorem}\label{t:main}
Qualia exist and are intrinsic to the experiences that have them, and the structure of each modality's qualia can be fully described relationally.
\end{theorem}

\begin{proof}
Existence is Proposition~\ref{p:exist}. A property of an experience either depends on the experience's relations to things that are not its constituents or does not (Definition~\ref{d:intrinsic}). The candidates for relational constitution are the quality's place among other qualities, the experience's representational relations, and its causal or functional relations. The first two are excluded by Lemmas~\ref{l:notrel} and \ref{l:notrep}. Functional relations are a species of relational structure, and the argument of Lemma~\ref{l:notrel} applies to them: a functional role is specified by a Ramsey sentence and is satisfied by any domain of realizers with the right causal structure. Hence qualia are intrinsic. The relational structure of the qualities of a modality is the structure $\langle Q, R_1,\ldots,R_n\rangle$ of Definition~\ref{d:space}, which a relational description specifies completely, since a relational description just is a specification of which tuples stand in which relations.
\end{proof}

\section{What relational description leaves out}

The two halves of Theorem~\ref{t:main} combine into a precise account of what can and cannot be said about qualia in structural terms.

\begin{proposition}[Completeness up to automorphism]\label{p:auto}
Let $D$ be a complete relational description of a quality space $\mathcal{Q} = \langle Q, R_1, \ldots, R_n\rangle$. Then $D$ determines the assignment of qualia to places exactly up to the automorphisms of $\mathcal{Q}$. If $\mathcal{Q}$ has no non-trivial automorphisms, $D$ determines each quale's place uniquely; it never determines which intrinsic character the occupant of a given place has.
\end{proposition}

\begin{proof}
Two assignments of qualia to places both satisfy $D$ if and only if they differ by a structure-preserving bijection of $\mathcal{Q}$ onto itself, that is, an automorphism (Lemma~\ref{l:newman}). If the only automorphism is the identity, the place of each quale is fixed: there is exactly one place that stands in the specified relations to all the others. But $D$ states only that some quality occupies each place; it does not say what that quality is like, since by Lemma~\ref{l:notrel} what it is like is not fixed by its relations.
\end{proof}

Proposition~\ref{p:auto} separates two questions that discussions of the inverted spectrum have run together. One is whether the structure of a quality space allows a non-trivial automorphism, so that two subjects could have the same relational structure while each place is occupied by a different quale for each. \citet{Palmer1999} argued that the asymmetries of human color space, such as the different numbers of discriminable shades between unique hues and the different lightness of saturated yellow and saturated blue, rule out most candidate inversions, leaving at most a few structure-preserving mappings. The other question is whether, given a fixed place, a relational description says what its occupant is like. Proposition~\ref{p:auto} answers the second question negatively whatever the answer to the first. MND-22 holds that spectrum inversion is possible but undetectable in principle; on the present account, any inversion that preserves all relations would be undetectable by any relational test, behavioral or neural, since all such tests probe structure.

The proposition also explains the ineffability that Lewis ascribed to qualia, without making them mysterious. Everything about a quale that can be communicated to someone who has not had it is relational: its similarities, its oppositions, its place in an order. What cannot be communicated is the character of the occupant, because communication of character, as opposed to structure, requires that the hearer be able to fix the reference of a term to the quality by having it. The companion manuscript MND-10-b develops this point. MND-03 applies the same distinction to echolocation: much of what it is like to be a bat can be specified structurally, as a sonar-structured spatial field, though not all. The unspecified residue is what Proposition~\ref{p:auto} predicts: the intrinsic characters of the places in a structure we can describe.

\section{Further cases}

\paragraph{Pain.} Pains vary in intensity, location, temporal profile and quality (burning, aching, stabbing), and these dimensions have a relational structure that pain scales and questionnaires capture. Theorem~\ref{t:main} says that the structure is fully describable and that the character of each pain is not exhausted by it. The second claim has the consequence that two pain systems with the same structure might differ in what their pains are like. It does not have the consequence that we cannot know whether another creature is in pain, since the evidence for pain in others is evidence of a structure of states and responses whose best explanation is that the creature has experiences of the relevant kind (see MND-06 and MND-23).

\paragraph{Afterimages.} An afterimage has a color quale and a location in the visual field but no object. Representationalists treat it as representing a colored patch that is not there. Whatever the merits of that analysis, the quale is intrinsic in the sense of Definition~\ref{d:intrinsic}: duplicating the afterimage experience duplicates its color.

\paragraph{Perceptual experience.} On the naïve realism defended in MND-16, the phenomenal character of a veridical perception is partly constituted by the perceived object and its properties. This is compatible with Theorem~\ref{t:main}, because the object is a constituent of the perceptual experience (Definition~\ref{d:intrinsic}), and a duplicate of the experience includes a duplicate of the object as perceived. What the theorem excludes is that the character of the experience be constituted by relations to things outside the experience, such as a functional role or a represented property that the experience does not contain.

\paragraph{Novel modalities.} A subject fitted with a sensory substitution device, or a creature with a modality we lack, may have a quality space whose relational structure we can describe fully. By Proposition~\ref{p:auto}, we can then know everything about its qualia except what they are like, which we cannot know without having them or qualities sufficiently like them.

\section{Recent literature}

Two strands of recent work bear directly on the result. The first is the illusionist programme. \citet{Frankish2016} distinguishes phenomenal properties from \enquote{quasi-phenomenal} properties, the physical properties that introspection misrepresents as phenomenal, and argues that the illusionist need posit no phenomenal seeming, only an introspective misrepresentation. This is the second horn of our dilemma, and we reply to it below. \citet{Balog2016} argues that illusionism cannot account for the phenomenal concepts with which we think about our experiences, and \citet{Schwitzgebel2016} that phenomenal consciousness can be defined by examples innocently enough that the illusionist must deny that there are any positive examples, which is implausible. ETH-07 relies on the existence of qualia, argued for here, in rejecting the illusionist objection to grounding moral status in sentience.

The second strand is the formal study of quality spaces. \citet{Lee2021} argues that mental qualities can be modeled as points in structured spaces and that such models capture much of what is distinctive about phenomenal character, including its unity and its differences between modalities. The present account takes such models to be complete relational descriptions in the sense of Definition~\ref{d:space}. Lee is neutral, as far as the models themselves go, on whether the qualities are exhausted by their places; Proposition~\ref{p:auto} gives a reason to deny that they are.

The view that results is congenial to Russellian monism. MND-01 argues that consciousness is the intrinsic nature of physical structure and that the explanatory gap arises from two modes of presentation of the same facts. On that view, physics describes the relational structure of the world and experience presents some intrinsic natures that realize it, which is the relation Theorem~\ref{t:main} finds between qualia and their structure. The theorem does not depend on Russellian monism, but it identifies the place where such a view must locate qualia.

\section{Objections and replies}

\begin{objection}[Quining]\label{o:quine}
Qualia are supposed to be ineffable, intrinsic, private and directly apprehensible, and nothing is all four \citep{Dennett1988}.
\end{objection}

\begin{reply}
Definition~\ref{d:quale} requires none of the four. Intrinsicness we argue for. Ineffability follows, in the limited form described after Proposition~\ref{p:auto}: qualia are not ineffable in the sense that nothing can be said about them, since their whole structure can be said, but in the sense that their character cannot be conveyed by description. Privacy and direct apprehension we do not claim, except to the extent of Principle~\ref{p:intro}, which allows error. Dennett's coffee tasters show that judgments about the causes and history of one's experience are unreliable, which MND-20 accepts; they do not show that there is no fact about how the coffee now tastes.
\end{reply}

\begin{objection}[Transparency]\label{o:transp}
When one attends to one's experience of a tomato, one finds only the redness of the tomato. There are no qualities of the experience to be found \citep{Harman1990}.
\end{objection}

\begin{reply}
Transparency, so far as it holds, concerns where one's attention goes in perception: to the object. It does not show that the experience has no qualitative properties, only that they are not presented as properties of something other than the object. On the naïve realist view the qualities one attends to in perception are qualities of the object as they figure in the experience, and the experience's phenomenal character includes them. In bodily sensation, afterimages and blurred vision, moreover, experience is not transparent: one can attend to the blur as a feature of seeing rather than of what is seen \citep{Kind2003}.
\end{reply}

\begin{objection}[Quasi-phenomenal properties]\label{o:frankish}
The illusionist does not posit a phenomenal seeming. She posits an introspective representation that misrepresents physical, quasi-phenomenal properties as phenomenal. Lemma~\ref{l:illusion} begs the question by assuming that introspection is reliable about the very matter in dispute \citep{Frankish2016}.
\end{objection}

\begin{reply}
Principle~\ref{p:intro} is not assumed for the case; it is argued for in MND-20 on general grounds, from the method by which judgments about current experience are formed and from the pattern of evidence for introspective error. The illusionist must show that this method fails in the one respect that would make illusionism true, and the evidence offered, which concerns confabulation, theory-laden classification and the unattended periphery, does not bear on it. Moreover, the illusionist must say what the misrepresentation represents. If it represents experiences as having properties in virtue of which there is something it is like to have them, then the illusionist owes an account of how we came by that concept, which on her view has no instances and no phenomenal seemings to be acquired from. \citet{Balog2016} presses this point.
\end{reply}

\begin{objection}[Epiphenomenalism]\label{o:epi}
If qualia are intrinsic and not constituted by causal role, they make no causal difference, and our judgments about them are not caused by them. Then the reliability of introspection about them is a coincidence.
\end{objection}

\begin{reply}
Intrinsic properties can be causally efficacious; the mass of a body is intrinsic to it and makes a difference to what it does. What Theorem~\ref{t:main} denies is that qualia are constituted by their causal roles, not that they have them. MND-11 argues that mental causes are proportional causes, so that the qualitative character of a pain, rather than the fine-grained physical state that realizes it, is the proportional cause of the judgment that one is in pain. Introspective reliability is then no coincidence.
\end{reply}

\begin{objection}[Structuralism]\label{o:struct}
Lemma~\ref{l:notrel} trades on the claim that a domain of numbers could be given the structure of color space. But no domain of numbers could stand in the relations of phenomenal similarity, which are themselves qualitative. The structuralist takes the relations, not the relata, to be phenomenal.
\end{objection}

\begin{reply}
If the relations of similarity are phenomenal relations that only qualia can stand in, then the relata must have some character that fits them for the relations, and that character is what the realist calls qualia. If the relations are characterized only by their formal properties, Newman's point applies. The structuralist must take one horn or the other, and on the first, structuralism is a description of the structure of intrinsic qualities, which is the view defended here.
\end{reply}

\section{What the result does not establish}

The result does not establish that qualia are non-physical. It is consistent with the physicalism of MND-02, with the phenomenal concept strategy of MND-04 and MND-05, and with the Russellian monism of MND-01, which is a form of physicalism on MND-02's definition. Nor does it say which creatures have qualia, which is the question of MND-23, or whether any artificial system does. It does not settle whether a given quality space has non-trivial automorphisms, which is an empirical question for each modality, or whether any inversion is actual.

\section{Conclusion}

Qualia exist: introspection of current experience, which is reliable, finds them, and illusionism can explain this datum away only by positing a seeming that is either phenomenal or not the datum. They are intrinsic: a relational description of a quality space, like any structural description, fixes its domain only up to isomorphism and cannot say what the qualities are like, and representational and functional relations fare no better. Their structure is nonetheless fully describable relationally, and what is left out of the description, by Proposition~\ref{p:auto}, is exactly the intrinsic character of each place.

\begin{thebibliography}{99}\small
\bibitem[Balog(2016)]{Balog2016} Balog, K. (2016). Illusionism's discontent. \emph{Journal of Consciousness Studies}, 23(11--12), 40--51.
\bibitem[Block(1990)]{Block1990} Block, N. (1990). Inverted Earth. \emph{Philosophical Perspectives}, 4, 53--79.
\bibitem[Block(1996)]{Block1996} Block, N. (1996). Mental paint and mental latex. \emph{Philosophical Issues}, 7, 19--49.
\bibitem[Chalmers(1996)]{Chalmers1996} Chalmers, D.~J. (1996). \emph{The Conscious Mind: In Search of a Fundamental Theory}. Oxford University Press.
\bibitem[Chalmers(2018)]{Chalmers2018} Chalmers, D.~J. (2018). The meta-problem of consciousness. \emph{Journal of Consciousness Studies}, 25(9--10), 6--61.
\bibitem[Clark(1993)]{Clark1993} Clark, A. (1993). \emph{Sensory Qualities}. Oxford University Press.
\bibitem[Dennett(1988)]{Dennett1988} Dennett, D.~C. (1988). Quining qualia. In A.~J. Marcel and E. Bisiach (eds.), \emph{Consciousness in Contemporary Science}, 42--77. Oxford University Press.
\bibitem[Dennett(1991)]{Dennett1991} Dennett, D.~C. (1991). \emph{Consciousness Explained}. Little, Brown.
\bibitem[Frankish(2016)]{Frankish2016} Frankish, K. (2016). Illusionism as a theory of consciousness. \emph{Journal of Consciousness Studies}, 23(11--12), 11--39.
\bibitem[Hardin(1988)]{Hardin1988} Hardin, C.~L. (1988). \emph{Color for Philosophers: Unweaving the Rainbow}. Hackett.
\bibitem[Harman(1990)]{Harman1990} Harman, G. (1990). The intrinsic quality of experience. \emph{Philosophical Perspectives}, 4, 31--52.
\bibitem[Jackson(1982)]{Jackson1982} Jackson, F. (1982). Epiphenomenal qualia. \emph{The Philosophical Quarterly}, 32(127), 127--136.
\bibitem[Kind(2003)]{Kind2003} Kind, A. (2003). What's so transparent about transparency? \emph{Philosophical Studies}, 115(3), 225--244.
\bibitem[Lee(2021)]{Lee2021} Lee, A.~Y. (2021). Modeling mental qualities. \emph{The Philosophical Review}, 130(2), 263--298.
\bibitem[Lewis(1929)]{Lewis1929} Lewis, C.~I. (1929). \emph{Mind and the World-Order}. Charles Scribner's Sons.
\bibitem[Lewis(1970)]{Lewis1970} Lewis, D. (1970). How to define theoretical terms. \emph{The Journal of Philosophy}, 67(13), 427--446.
\bibitem[Nagel(1974)]{Nagel1974} Nagel, T. (1974). What is it like to be a bat? \emph{The Philosophical Review}, 83(4), 435--450.
\bibitem[Newman(1928)]{Newman1928} Newman, M.~H.~A. (1928). Mr.\ Russell's \enquote{causal theory of perception}. \emph{Mind}, 37(146), 137--148.
\bibitem[Palmer(1999)]{Palmer1999} Palmer, S.~E. (1999). Color, consciousness, and the isomorphism constraint. \emph{Behavioral and Brain Sciences}, 22(6), 923--943.
\bibitem[Schwitzgebel(2016)]{Schwitzgebel2016} Schwitzgebel, E. (2016). Phenomenal consciousness, defined and defended as innocently as I can manage. \emph{Journal of Consciousness Studies}, 23(11--12), 224--235.
\bibitem[Shoemaker(1982)]{Shoemaker1982} Shoemaker, S. (1982). The inverted spectrum. \emph{The Journal of Philosophy}, 79(7), 357--381.
\bibitem[Tye(1995)]{Tye1995} Tye, M. (1995). \emph{Ten Problems of Consciousness: A Representational Theory of the Phenomenal Mind}. MIT Press.
\end{thebibliography}

\end{document}
